\section{Security analysis}\label{sec:security}

SA-HEACS removes code recovery (G1) for cloud-only features, because the code is not on the
device; it reduces but does not remove functional cloning (G2) and free-riding (G3). The
original blueprint rated every row \enquote{fully mitigated}; \cref{tab:threats} separates
what the architecture eliminates from what it only makes harder.

\begin{table}[htbp]
\centering
\caption{Reverse-engineering threats under client-only protection and under SA-HEACS.}
\label{tab:threats}
\small
\begin{tabularx}{\linewidth}{@{}>{\raggedright\arraybackslash}p{3.1cm} Y Y Y@{}}
\toprule
Threat & Client-only protection & SA-HEACS, cloud-only features & Why \\
\midrule
Static decompilation (Ghidra, IDA)
  & Delayed by obfuscation \cite{collberg1997taxonomy,schrittwieser2016obfuscation}
  & Eliminated for the protected code
  & Code is not in the client package \\
Memory dumping and hooking (Frida)
  & Delayed by anti-debug and integrity checks
  & Eliminated for the protected code; client inputs and outputs remain visible
  & Computation runs server-side \\
AI-assisted decompilation and summarisation
  & Increasingly effective \cite{chen2022dirty,tan2024llm4decompile}
  & Recovers only the shell and the API schema
  & Nothing protected to lift \\
Black-box functional cloning via the API
  & Possible
  & Remains possible; slowed by rate limits, monitoring, output shaping
  & Outputs leak function \cite{tramer2016stealing,jagielski2020extraction,orekondy2019knockoff} \\
Patched or cloned client driving the API
  & Possible
  & Raised cost: attestation, per-account quotas
  & Attestation can be bypassed on rooted or emulated devices \\
Shell and UI cloning
  & Possible & Accepted & Clone has no protected capability \\
Cost exhaustion (DoS on GPUs)
  & Not applicable & New risk & Remote compute is costly per request \\
\bottomrule
\end{tabularx}
\end{table}

\subsection{What is eliminated}
Against G1, the argument is simple and does not depend on obfuscation strength: an analyst
cannot recover an artefact that never reaches their machine. Barak et al.'s impossibility
result \cite{barak2012impossibility} limits what any shipped program can hide; SA-HEACS
avoids that limit by not shipping the program. This holds only if the server's responses do
not contain the code or weights in another form, so responses must not include debugging
traces, intermediate tensors or model files.

\subsection{What remains: extraction through the API}
A server-side function is still an oracle. Tramèr et al.\ \cite{tramer2016stealing} and
Orekondy et al.\ \cite{orekondy2019knockoff} show that query access alone can produce a
useful substitute model, and an AI agent can automate the query campaign as easily as it
automates decompilation. SA-HEACS therefore treats G2 as a residual risk with layered
controls:
\begin{itemize}
  \item \textbf{Per-account and per-device quotas} tied to attested identities, so that
        query volume costs real accounts and devices.
  \item \textbf{Query-distribution monitoring} in the style of PRADA \cite{juuti2019prada},
        flagging accounts whose inputs look synthetic or adversarially spaced.
  \item \textbf{Output shaping}: return the product-level result (for example an 8-bit
        image) rather than logits, confidences or high-precision intermediates that speed up
        extraction \cite{tramer2016stealing}.
  \item \textbf{Watermarking} of model behaviour or outputs \cite{adi2018watermarking}, so
        that a stolen substitute can be identified even if it cannot be prevented.
\end{itemize}
These controls raise the price of a clone; none makes it impossible. The honest claim is
that SA-HEACS converts a white-box attack into a black-box one and makes the black-box attack
measurable and rate-limited.

\subsection{Free-riding and attestation}
Attestation \cite{rfc9334,google_playintegrity,apple_appattest} lets the server refuse
requests from modified clients, but attackers can bypass it on compromised devices, and a
legitimate unmodified client can still be scripted through accessibility or UI automation.
Attestation is therefore one signal for the rate limiter and anomaly detector, not a gate
that is assumed unbreakable.

\subsection{New attack surface}
Moving IP to the server creates risks that client-only software does not have:
cost-exhaustion attacks against GPU capacity, injection or deserialisation bugs in the RPC
layer, and cross-tenant leakage if batching or caching mixes state between users.
Micro-batching must keep per-row isolation, and prefetch caches must be keyed per user.
